\providecommand{\Needspace}[1]{}
\Needspace{8\baselineskip}
\item \textbf{Limpieza y resumen de observaciones astronómicas con pandas}.
\nopagebreak[4]

\begin{tcolorbox}[colback=blue!2!white,colframe=blue!55!black,title={Enunciado},
                  before skip=3pt,after skip=5pt,boxsep=3pt,top=3pt,bottom=3pt,
                  breakable]
\small
\raggedright
Se tiene un conjunto pequeño de observaciones astronómicas. Cada observación
tiene un nombre, un tipo de objeto, un flujo medido, una distancia y una
estimación del ruido instrumental.

Use el siguiente bloque para crear el DataFrame:

\begin{pythonjupyter}
datos = {
    "nombre": ["A1", "A2", "B1", "B2", "C1", "C2", "D1", "D2"],
    "tipo": ["estrella", "estrella", "galaxia", "galaxia",
             "cuásar", "cuásar", "estrella", "galaxia"],
    "flujo": [12.5, 8.1, 20.4, None, 5.2, 6.3, 10.9, 18.7],
    "distancia_pc": [4.2, 7.8, 120.0, 98.0, 300.0, None, 5.5, 110.0],
    "ruido": [0.8, 0.9, 2.5, 2.1, 0.6, 0.7, None, 2.2],
}
\end{pythonjupyter}

Realice las siguientes operaciones:

\begin{itemize}
\setlength{\topsep}{4pt}
\setlength{\itemsep}{3pt}
\setlength{\parsep}{0pt}
\item[(a)] Importe \texttt{pandas}, \texttt{numpy} y
\texttt{matplotlib.pyplot}. Cree un DataFrame llamado \texttt{obs} a partir del
diccionario \texttt{datos}.
\item[(b)] Inspeccione el DataFrame mostrando las primeras filas, dimensiones,
tipos de datos y cantidad de valores faltantes por columna.
\item[(c)] Cree un DataFrame llamado \texttt{obs\_limpias} eliminando las filas
con datos faltantes en \texttt{flujo}, \texttt{distancia\_pc} o \texttt{ruido}.
Use \texttt{copy()}.
\item[(d)] En \texttt{obs\_limpias}, cree una columna llamada \texttt{snr} usando
\[
\texttt{snr}=\frac{\texttt{flujo}}{\texttt{ruido}}.
\]
\item[(e)] Agrupe \texttt{obs\_limpias} por \texttt{tipo}. Calcule la cantidad
de observaciones por tipo usando la columna \texttt{tipo}.
\item[(f)] Cree una copia sin la columna \texttt{nombre}. Agrupe esa copia por
\texttt{tipo} y calcule el promedio y la desviación estándar por tipo. Use
\texttt{numeric\_only=True} si lo necesita.
\item[(g)] Filtre las observaciones con \texttt{snr > 10} y
\texttt{distancia\_pc < 150}. Muestre solamente \texttt{nombre}, \texttt{tipo},
\texttt{flujo}, \texttt{distancia\_pc} y \texttt{snr}.
\item[(h)] Haga un gráfico de dispersión simple con \texttt{distancia\_pc} en el
eje horizontal y \texttt{snr} en el eje vertical. Agregue etiquetas, título y
grilla.
\end{itemize}
\end{tcolorbox}

\begin{solution}
\begin{pythonjupyter}
import pandas as pd
import numpy as np
import matplotlib.pyplot as plt

datos = {
    "nombre": ["A1", "A2", "B1", "B2", "C1", "C2", "D1", "D2"],
    "tipo": ["estrella", "estrella", "galaxia", "galaxia",
             "cuásar", "cuásar", "estrella", "galaxia"],
    "flujo": [12.5, 8.1, 20.4, None, 5.2, 6.3, 10.9, 18.7],
    "distancia_pc": [4.2, 7.8, 120.0, 98.0, 300.0, None, 5.5, 110.0],
    "ruido": [0.8, 0.9, 2.5, 2.1, 0.6, 0.7, None, 2.2],
}

obs = pd.DataFrame(datos)

print(obs.head())
print("Dimensiones:", obs.shape)
print(obs.dtypes)
print("Valores faltantes:")
print(obs.isna().sum())

columnas_necesarias = ["flujo", "distancia_pc", "ruido"]
obs_limpias = obs.dropna(subset=columnas_necesarias).copy()

obs_limpias["snr"] = obs_limpias["flujo"] / obs_limpias["ruido"]

grupo_tipo = obs_limpias.groupby("tipo")
cantidad_por_tipo = grupo_tipo["tipo"].count()
print(cantidad_por_tipo)

obs_sin_nombre = obs_limpias.drop(columns=["nombre"]).copy()
grupo_tipo_num = obs_sin_nombre.groupby("tipo")

promedio_por_tipo = grupo_tipo_num.mean(numeric_only=True)
desviacion_por_tipo = grupo_tipo_num.std(numeric_only=True)

print(promedio_por_tipo)
print(desviacion_por_tipo)

filtro = (
    (obs_limpias["snr"] > 10) &
    (obs_limpias["distancia_pc"] < 150)
)

seleccion = obs_limpias.loc[
    filtro,
    ["nombre", "tipo", "flujo", "distancia_pc", "snr"]
]
print(seleccion)

plt.figure(figsize=(7, 5))
plt.scatter(obs_limpias["distancia_pc"], obs_limpias["snr"])
plt.xlabel("Distancia [pc]")
plt.ylabel("SNR")
plt.title("SNR versus distancia")
plt.grid(True, alpha=0.3)
plt.show()
\end{pythonjupyter}
\end{solution}
